\documentclass[11pt, a4paper]{article}

% bfw-style.tex — gemeinsame Style-Definitionen für alle Handouts/Aufgabenhefte
% Voraussetzung: Kompilation mit xelatex (wegen fontspec)

\usepackage[ngerman]{babel}
\usepackage{geometry}
\geometry{a4paper, top=2.2cm, bottom=2.2cm, left=2.3cm, right=2.3cm}

\usepackage{fontspec}
% Fallback-Kette: Source Sans Pro -> Calibri -> Segoe UI -> Latin Modern Sans
% (Latin Modern ist immer mit MiKTeX vorhanden)
\IfFontExistsTF{Source Sans Pro}{%
  \setmainfont{Source Sans Pro}[Ligatures=TeX]%
}{\IfFontExistsTF{Calibri}{%
  \setmainfont{Calibri}[Ligatures=TeX]%
}{\IfFontExistsTF{Segoe UI}{%
  \setmainfont{Segoe UI}[Ligatures=TeX]%
}{%
  \setmainfont{Latin Modern Sans}[Ligatures=TeX]%
}}}
\IfFontExistsTF{Source Code Pro}{%
  \setmonofont{Source Code Pro}[Scale=0.92]%
}{\IfFontExistsTF{Consolas}{%
  \setmonofont{Consolas}[Scale=0.92]%
}{%
  \setmonofont{Latin Modern Mono}[Scale=0.92]%
}}

\usepackage{xcolor}
% Farbpalette: hell, freundlich, modern
\definecolor{bfwPrimary}{HTML}{0EA5A4}    % Petrol/Türkis - Primär
\definecolor{bfwPrimaryDark}{HTML}{0F766E} % dunkler Petrol für Headings
\definecolor{bfwAccent}{HTML}{F59E0B}     % Amber - Akzent
\definecolor{bfwSuccess}{HTML}{10B981}    % Grün - Erfolg / Lösung
\definecolor{bfwWarn}{HTML}{F97316}       % Orange - Achtung
\definecolor{bfwInk}{HTML}{0F172A}        % fast Schwarz
\definecolor{bfwInkSoft}{HTML}{475569}    % gedämpftes Grau
\definecolor{bfwBgSoft}{HTML}{F8FAFC}     % off-white
\definecolor{bfwBgTeal}{HTML}{ECFEFF}     % helles Türkis
\definecolor{bfwBgAmber}{HTML}{FEF3C7}    % helles Gelb
\definecolor{bfwBgRose}{HTML}{FFE4E6}     % helles Rosé für Eselsbrücken

\usepackage[colorlinks=true, linkcolor=bfwPrimaryDark, urlcolor=bfwPrimaryDark]{hyperref}
\usepackage{enumitem}
\setlist[itemize]{leftmargin=*, itemsep=2pt, topsep=2pt}
\setlist[enumerate]{leftmargin=*, itemsep=2pt, topsep=2pt}

\usepackage[most]{tcolorbox}
\usepackage{booktabs}
\usepackage{tikz}
\usetikzlibrary{decorations.pathmorphing, positioning, shapes.geometric, arrows.meta}

% Merkkästchen — wichtige Definition / Kernpunkt
\newtcolorbox{merksatz}[1][]{%
  enhanced, breakable,
  colback=bfwBgTeal,
  colframe=bfwPrimary,
  boxrule=0.6pt,
  arc=4pt,
  left=10pt, right=10pt, top=8pt, bottom=8pt,
  fonttitle=\bfseries\color{bfwPrimaryDark},
  title={\faLightbulb\ Merksatz},
  #1
}

% Eselsbrücke — Mnemoniks
\newtcolorbox{eselsbruecke}[1][]{%
  enhanced, breakable,
  colback=bfwBgRose,
  colframe=bfwAccent,
  boxrule=0.6pt,
  arc=4pt,
  left=10pt, right=10pt, top=8pt, bottom=8pt,
  fonttitle=\bfseries\color{bfwWarn},
  title={Eselsbrücke},
  #1
}

% Beispiel-Box
\newtcolorbox{beispiel}[1][]{%
  enhanced, breakable,
  colback=bfwBgSoft,
  colframe=bfwInkSoft,
  boxrule=0.4pt,
  arc=3pt,
  left=10pt, right=10pt, top=8pt, bottom=8pt,
  fonttitle=\bfseries\color{bfwInk},
  title={Beispiel},
  #1
}

% Lösungs-Box (für Aufgabenhefte)
\newtcolorbox{loesung}[1][]{%
  enhanced, breakable,
  colback=bfwBgAmber!40,
  colframe=bfwSuccess,
  boxrule=0.6pt,
  arc=3pt,
  left=10pt, right=10pt, top=8pt, bottom=8pt,
  fonttitle=\bfseries\color{bfwSuccess!70!black},
  title={Lösung},
  #1
}

% Glossar-Eintrag
\newcommand{\glossareintrag}[2]{%
  \par\noindent\textbf{\color{bfwPrimaryDark}#1} \enspace #2\par\smallskip
}

% Section-Stil — etwas farbiger als Standard
\usepackage{titlesec}
\titleformat{\section}
  {\Large\bfseries\color{bfwPrimaryDark}}
  {\thesection}{0.6em}{}
\titleformat{\subsection}
  {\large\bfseries\color{bfwPrimary}}
  {\thesubsection}{0.5em}{}
\titleformat{\subsubsection}
  {\normalsize\bfseries\color{bfwInk}}
  {\thesubsubsection}{0.4em}{}

% TikZ-Stil "skizzenhaft" — etwas weicher als Rough.js, aber stabil
\tikzset{
  roughbox/.style = {
    draw=bfwPrimaryDark, thick, fill=bfwBgTeal,
    rounded corners=4pt, inner sep=6pt
  },
  roughaccent/.style = {
    draw=bfwAccent, thick, fill=bfwBgAmber,
    rounded corners=4pt, inner sep=6pt
  },
  roughline/.style = {
    draw=bfwInkSoft, thick, -{Stealth[length=2.5mm]}
  }
}

% Bullet-Symbole (bewusst kein FontAwesome — vermeidet Font-Installations-Probleme)
\newcommand{\faLightbulb}{$\bullet$}
\newcommand{\faBookmark}{$\bullet$}

% Header/Footer
\usepackage{fancyhdr}
\pagestyle{fancy}
\fancyhf{}
\renewcommand{\headrulewidth}{0pt}
\fancyfoot[L]{\small\color{bfwInkSoft}\thepage}
\fancyfoot[R]{\small\color{bfwInkSoft} BFW M-064 Beschaffung im Büromanagement}


\title{\vspace*{-1.5cm}\Huge\bfseries\color{bfwPrimaryDark}Aufgabenheft Tag~7\\[0.4em]
  \large\color{bfwInkSoft}\textnormal{Eingangsrechnungen --- Übungen im IHK-Klausurformat}}
\author{\textsf{Beschaffung im Büromanagement}\\
\textsf{\small\copyright\ Can Siebert\,\textperiodcentered\,2026}}
\date{Selbstlernmaterial \textperiodcentered{} 20.05.2026}

\begin{document}
\maketitle
\thispagestyle{empty}

\begin{merksatz}[title={\bfseries So nutzen Sie dieses Aufgabenheft}]
Bearbeiten Sie alle Aufgaben zunächst \emph{ohne} in die Lösungen zu schauen.
Beachten Sie die \emph{Operatoren} und schreiben Sie Antworten in vollständigen Sätzen.
Lösungen ab Seite~\pageref{loesungen}.
\end{merksatz}

\tableofcontents
\vspace{1em}
\hrule

\section{Aufgaben}

\subsection*{Aufgabe~1 --- Erfüllung und Eigentumserwerb (10 Punkte)}

\begin{enumerate}[label=(\alph*)]
  \item \textbf{Erläutern} Sie, wann ein Kaufvertrag nach §\,362 BGB erfüllt ist. \hfill (2~P.)
  \item \textbf{Nennen} Sie die Voraussetzungen des \emph{gutgläubigen} Eigentumserwerbs nach §\,932 BGB. \hfill (4~P.)
  \item \textbf{Beurteilen} Sie folgenden Fall: Sie kaufen vom Großhandel \emph{Office-Supply GmbH} einen gebrauchten Drucker. Später stellt sich heraus, der Verkäufer war nicht Eigentümer. Können Sie trotzdem Eigentum erworben haben? Begründen Sie. \hfill (4~P.)
\end{enumerate}

\subsection*{Aufgabe~2 --- Eigentumsvorbehalt (12 Punkte)}

\noindent\textbf{\color{bfwAccent}YouTube-Vorbereitung:}\enspace \href{https://www.youtube.com/watch?v=vRbdTNv1Xys}{Eigentumsvorbehalt einfach / verlängert / erweitert} (\url{https://www.youtube.com/watch?v=vRbdTNv1Xys})

\medskip
\begin{enumerate}[label=(\alph*)]
  \item \textbf{Erläutern} Sie den Begriff \emph{Eigentumsvorbehalt} nach §\,449 BGB. \hfill (2~P.)
  \item \textbf{Unterscheiden} Sie die vier Formen des Eigentumsvorbehalts (einfach, erweitert, verlängert, weitergeleitet). \hfill (8~P.)
  \item \textbf{Beurteilen} Sie: Welche Form ist für einen Lieferanten von Tonerkartuschen, dessen Käufer Wiederverkäufer ist, am sinnvollsten? Begründen Sie. \hfill (2~P.)
\end{enumerate}

\subsection*{Aufgabe~3 --- Erfüllungsort und Gefahrübergang (10 Punkte)}

\begin{enumerate}[label=(\alph*)]
  \item \textbf{Unterscheiden} Sie Holschuld, Bringschuld und Schickschuld. \hfill (3~P.)
  \item \textbf{Erläutern} Sie §\,447 BGB --- wann geht beim Versendungskauf die Gefahr auf den Käufer über? \hfill (3~P.)
  \item \textbf{Beurteilen} Sie: Die \emph{Müller GmbH} bestellt 50~Bürostühle zu sich ins Haus. Der Spediteur übernimmt die Stühle, sie werden auf dem Transport beschädigt. Wer trägt die Gefahr (B2B Versendungskauf)? \hfill (4~P.)
\end{enumerate}

\subsection*{Aufgabe~4 --- Pflichtangaben einer Rechnung (12 Punkte)}

\begin{enumerate}[label=(\alph*)]
  \item \textbf{Nennen} Sie acht Pflichtangaben einer Rechnung nach §\,14 UStG. \hfill (8~P.)
  \item \textbf{Erläutern} Sie, warum die Vollständigkeit dieser Angaben für den \emph{Vorsteuerabzug} entscheidend ist. \hfill (2~P.)
  \item \textbf{Beurteilen} Sie: Welche Folge hat eine fehlende Steuernummer auf der Eingangsrechnung? \hfill (2~P.)
\end{enumerate}

\subsection*{Aufgabe~5 --- Rechnungsprüfung in der Praxis (10 Punkte)}

Sie erhalten als Sachbearbeiterin der \emph{Schreibwaren-Vertrieb GmbH} folgende
Rechnung:

\begin{itemize}[itemsep=0pt]
  \item 50~Stück Aktenordner Standard à 4,80~€ netto
  \item 10~Pakete Druckerpapier à 28,00~€ netto
  \item Versand: 12,00~€ netto
  \item USt 19\,\%
\end{itemize}

\begin{enumerate}[label=(\alph*)]
  \item \textbf{Berechnen} Sie die Zwischensumme netto, die Umsatzsteuer und den Bruttobetrag. \hfill (4~P.)
  \item \textbf{Erläutern} Sie den Drei-Punkt-Vergleich (Bestellung --- Wareneingang --- Rechnung). \hfill (3~P.)
  \item \textbf{Unterscheiden} Sie formelle, sachliche und rechnerische Prüfung. \hfill (3~P.)
\end{enumerate}

\subsection*{Aufgabe~6 --- IHK-Wissensfragen (6 Punkte)}

\begin{enumerate}[label=(\alph*)]
  \item \textbf{Erläutern} Sie kurz: Was ist die Folge, wenn eine Eingangsrechnung formelle Mängel aufweist? \hfill (2~P.)
  \item \textbf{Beurteilen} Sie: Warum darf der Lieferant beim verlängerten Eigentumsvorbehalt erlauben, dass der Käufer die Ware weiterverkauft? \hfill (2~P.)
  \item \textbf{Nennen} Sie zwei typische Fehler, die bei der Rechnungsprüfung auffallen können. \hfill (2~P.)
\end{enumerate}

\vspace{1em}
\noindent\textbf{Punkte gesamt: 60}

\newpage

\section{Lösungen}\label{loesungen}

\subsection*{Lösung Aufgabe~1}

\begin{loesung}
\textbf{(a)} Ein Kaufvertrag ist nach §\,362 BGB erfüllt, wenn die geschuldete Leistung
ordnungsgemäß bewirkt wird --- der Verkäufer übergibt die Sache und verschafft das
Eigentum, der Käufer zahlt den Kaufpreis und nimmt die Sache ab.

\textbf{(b)} Voraussetzungen § 932 BGB: (1)~Verkehrsgeschäft, (2)~bewegliche Sache,
(3)~Übergabe (oder Übergabesurrogat), (4)~guter Glaube an das Eigentum des
Veräußerers --- nicht grob fahrlässige Unkenntnis.

\textbf{(c)} Ja, Sie können nach §\,932 BGB Eigentum erworben haben, wenn Sie weder
wussten noch grob fahrlässig nicht wussten, dass der Verkäufer nicht Eigentümer war.
Voraussetzung: kein Fall des §\,935 BGB (z.\,B.\ Drucker war nicht gestohlen).
\end{loesung}

\subsection*{Lösung Aufgabe~2}

\begin{loesung}
\textbf{(a)} Eigentumsvorbehalt nach §\,449 BGB: Klausel im Kaufvertrag, durch die
das Eigentum erst mit voller Bezahlung des Kaufpreises auf den Käufer übergeht.

\textbf{(b)} Vier Formen:
\begin{itemize}[itemsep=0pt]
  \item \emph{Einfacher EV}: Eigentum geht mit Zahlung der konkreten Lieferung über.
  \item \emph{Erweiterter EV}: Eigentum geht erst über, wenn alle Forderungen aus der Geschäftsverbindung beglichen sind (Kontokorrent-Vorbehalt).
  \item \emph{Verlängerter EV}: Käufer darf weiterverkaufen; Forderungen aus dem Weiterverkauf werden im Voraus an den Lieferanten abgetreten (§ 398 BGB).
  \item \emph{Weitergeleiteter EV}: Vorbehaltsrecht wird auf einen Dritten übertragen (z.\,B.\ bei Verarbeitung in eine Maschine eines anderen Eigentümers).
\end{itemize}

\textbf{(c)} \emph{Verlängerter Eigentumsvorbehalt}: Der Wiederverkäufer braucht die
Erlaubnis, die Ware weiterzuverkaufen --- gleichzeitig sichert sich der Lieferant
durch die Vorausabtretung der Kundenforderungen ab.
\end{loesung}

\subsection*{Lösung Aufgabe~3}

\begin{loesung}
\textbf{(a)}
\begin{itemize}[itemsep=0pt]
  \item \emph{Holschuld}: Erfüllungsort beim Schuldner --- Käufer holt die Ware ab. (Gesetzlicher Regelfall, § 269 BGB.)
  \item \emph{Bringschuld}: Erfüllungsort beim Gläubiger --- Verkäufer bringt die Ware hin.
  \item \emph{Schickschuld}: Erfüllungsort beim Schuldner, aber Versand auf Käufergefahr.
\end{itemize}

\textbf{(b)} § 447 BGB: Beim Versendungskauf geht die Preisgefahr bereits mit
\emph{Übergabe an den Spediteur, Frachtführer oder die zur Ausführung der Versendung
bestimmte Person} auf den Käufer über --- nicht erst mit Ankunft.

\textbf{(c)} Im B2B-Versendungskauf trägt der Käufer (Müller GmbH) ab Übergabe an den
Spediteur die Gefahr (§ 447 BGB). Da die Stühle erst auf dem Transport beschädigt
wurden, geht der Schaden zu Lasten der Müller GmbH --- sie muss zahlen, kann aber
ggf.\ den Frachtführer in Anspruch nehmen. Wäre die Müller GmbH ein Verbraucher,
wäre § 475 II BGB anwendbar (Gefahrübergang erst bei Übergabe an Verbraucher).
\end{loesung}

\subsection*{Lösung Aufgabe~4}

\begin{loesung}
\textbf{(a)} Acht Pflichtangaben (§ 14 IV UStG):
\begin{enumerate}[itemsep=0pt]
  \item Vollständiger Name und Anschrift Lieferant + Empfänger.
  \item Steuernummer oder USt-IdNr.\ des Lieferanten.
  \item Ausstellungsdatum.
  \item Fortlaufende Rechnungsnummer.
  \item Menge und handelsübliche Bezeichnung der Ware/Leistung.
  \item Liefer- oder Leistungsdatum.
  \item Nettoentgelt nach Steuersätzen aufgeschlüsselt.
  \item Steuersatz und Steuerbetrag (oder Hinweis auf Steuerbefreiung).
\end{enumerate}

\textbf{(b)} Der Vorsteuerabzug nach § 15 UStG setzt eine \emph{ordnungsgemäße}
Rechnung voraus. Fehlt eine Pflichtangabe, kann das Finanzamt den Vorsteuerabzug
versagen --- die Umsatzsteuer wird zur echten Belastung des Empfängers.

\textbf{(c)} Folge fehlende Steuernummer: Die Rechnung ist formell mangelhaft, der
Vorsteuerabzug ist gefährdet. Der Empfänger sollte die Rechnung an den Lieferanten
zur Korrektur zurücksenden, bevor er zahlt.
\end{loesung}

\subsection*{Lösung Aufgabe~5}

\begin{loesung}
\textbf{(a)}
\begin{itemize}[itemsep=0pt]
  \item Aktenordner: $50 \times 4{,}80 = 240{,}00$~€
  \item Druckerpapier: $10 \times 28{,}00 = 280{,}00$~€
  \item Versand: $12{,}00$~€
  \item \textbf{Zwischensumme netto: $\mathbf{532{,}00~€}$}
  \item USt 19\,\%: $532{,}00 \times 0{,}19 = \mathbf{101{,}08~€}$
  \item \textbf{Bruttobetrag: $\mathbf{633{,}08~€}$}
\end{itemize}

\textbf{(b)} Drei-Punkt-Vergleich: Die Rechnung wird mit \emph{Bestellung}
(Was war vereinbart?) und \emph{Lieferschein/Wareneingang} (Was ist geliefert worden?)
abgeglichen. Stimmen alle drei überein, ist die Rechnung sachlich korrekt.

\textbf{(c)} Drei Stufen:
\begin{itemize}[itemsep=0pt]
  \item \emph{Formell}: Pflichtangaben § 14 UStG vorhanden?
  \item \emph{Sachlich}: Stimmen Lieferung/Leistung und Rechnung mit der Bestellung überein?
  \item \emph{Rechnerisch}: Sind Mengen $\times$ Preise, USt und Brutto richtig berechnet?
\end{itemize}
\end{loesung}

\subsection*{Lösung Aufgabe~6}

\begin{loesung}
\textbf{(a)} Bei formellen Mängeln (z.\,B.\ fehlende Pflichtangabe) ist der
Vorsteuerabzug gefährdet. Die Rechnung sollte zur Berichtigung an den Lieferanten
zurückgesandt werden.

\textbf{(b)} Beim verlängerten Eigentumsvorbehalt ermöglicht der Lieferant dem
Käufer den Weiterverkauf, weil der Käufer (z.\,B.\ Wiederverkäufer) sonst sein
Geschäftsmodell nicht betreiben könnte. Der Lieferant ist abgesichert, weil die
Forderungen aus dem Weiterverkauf bereits im Voraus an ihn abgetreten sind
(§ 398 BGB).

\textbf{(c)} Typische Fehler in Eingangsrechnungen:
\begin{itemize}[itemsep=0pt]
  \item Falscher Steuersatz (z.\,B.\ 7\,\% statt 19\,\%)
  \item Falsche Mengen oder Einzelpreise (Abweichung von der Bestellung)
  \item Fehlende Pflichtangaben (Steuernummer, Rechnungsnummer)
  \item Doppelberechnung von Versand- oder Verpackungskosten
\end{itemize}
\end{loesung}

\vfill
\hrule
\vspace{0.5em}
\noindent\small\color{bfwInkSoft}\copyright\ Can Siebert\,\textperiodcentered\,2026\,\textperiodcentered{} Beschaffung im Büromanagement\,\textperiodcentered{} Tag~7 Aufgabenheft

\end{document}
